\section{Introduction}
Researchers often reuse one predicted item label across many records and
construct several aggregates from those records. A website's content label
can enter the browsing histories of many people; a merchant classification
can enter many households' expenditures; a document label can contribute to
several measures of an organization or period. The resulting totals and
shares may become regression outcomes, covariates, or both. A mistake made
once at the item level then appears repeatedly in the analysis.

The consequences depend on the reported quantity. An error on an item used
widely by two groups can move both groups' means while largely canceling from
their difference. The same error can change several regression covariates
together. If the constructed measure is an outcome, a least-squares
coefficient is linear in its item labels. If it is a covariate, changing a
label changes the regression's numerator and denominator. Neither the number
of mistakes nor the number of downstream records describes these effects by
itself.

The distinction between accurate classification and accurate aggregate
estimation is established. Quantification addresses class prevalence;
cost quantification already allows nonnegative item weights
\citep{forman2008,hopkins2010,vaz2019quantification}. Research on predicted
covariates and outcomes shows that substituting predictions can change
regression coefficients and inference
\citep{fong2021predictions,wang2020postprediction}.
Classification followed by aggregation into regression indices is also
covered by recent generated-variable work
\citep{battaglia2025generated,christensen2026guide}. The present question is
how those inferential claims change when an uncertain item supplies parts of
many regression observations and several constructed variables.

Two other literatures clarify this question. Generated-regressor and
measurement-error methods account for estimated calibration, correlated
errors, and uncertainty in multiple estimation stages
\citep{pagan1984,murphy1985twostep,shaw2021correlated,boe2024sandwich}.
Shared-source inference explains why records with overlapping loadings need
not provide independent information
\citep{adao2019,borusyak2022shiftshare,owen2012arrays}. These results require
particular centering, sampling, or shock-assignment assumptions. A common
covariance structure alone does not establish those assumptions for
classification errors. Conversely, a biased coefficient does not imply an
inconsistent standard error: generated-covariate models include regimes in
which conventional robust SEs estimate variability correctly but ignore
first-order centering bias \citep{christensen2026guide,battaglia2025generated}.

I formalize reuse by matrices mapping item labels into record-level measures.
This representation yields the error in a linear estimate, the joint
covariance across estimates, and the normal-equation distortion when an
aggregate is a regression covariate. It also identifies the validation unit.
Verifying one item can correct pieces of many records without revealing any
record's complete covariate. Existing validation-sample regression methods,
survey difference estimators, and prediction-powered inference provide
correction principles
\citep{shepherd2012audit,lumley2011calibration,angelopoulos2023ppi,egami2023,
kluger2025imputed}. Applying them here requires preserving that item-level
construction, including label pairs in quadratic regression moments.

The baseline target is the coefficient that complete labels would produce
in the observed panel. For an existing probability sample of verified items,
I give moment corrections and an estimable joint covariance, then state the
conditions needed to propagate them into regression SEs. The corrections use
standard finite-population methods, including the treatment of quadratic
totals \citep{sitter2002quadratic}. Their exact unbiasedness does not guarantee
normal coefficient intervals. I also show what is lost when verified errors
are reduced to confusion counts and distinguish these fixed-label results
from population interpretations based on calibrated conditional means.

Simulations compare coefficient bias, empirical variability, estimated SEs,
and interval coverage for aggregates used on either or both regression sides.
A June 2022 YouGov browsing sample supplies an empirical illustration of
reuse and visit concentration. Its imposed errors are controlled examples,
not estimates of a fitted classifier's performance. A separate comparison of
observed adult-domain classifications shows how actual label disagreement
changes weighted shares without treating either classifier as verified truth.
